\section{Discussion}
\label{sec:discussion}

\subsection{Interpreting Weight Transfer Stability Under Cross-Language Evaluation}
\label{subsec:weight_stability_discussion}

Weights determined by optimization in Phase~1 ($w_{BL}=0.35$, $w_{WL}=0.40$, $w_{CES}=0.25$) yield 87.25\% accuracy using a dataset composed of C/C++/Java without any parameter fine-tuning. The result is consistent with, but cannot be seen as a validation of, universal generalizability of weights across languages. The setting is a perfect one where the system would be expected to be stable: All the submitted programs belong to the class of CS-1 imperative programs; both C and C++ use the same CPG frontend extractor; the weights were grid-searched on the same 20-problem dataset used for the Phase~2 experiments; and the structure of the program space is strictly restricted to loops, conditionals, and simple recursion.

To test the null hypothesis of equivalent performance, the \textbf{two-proportion z-test}~\cite{ref_fleiss2003} was performed on the C++ monolingual cohort ($360/400$ correct, $\hat{p}_1 = 0.900$) versus the cross-language cohort ($349/400$ correct, $\hat{p}_2 = 0.8725$). A two-proportion z-test, and not a McNemar test, is more appropriate in our case because these are distinct samples of 400 student programs (C++ monolingual cohort) versus 130 C, 130 C++, and 140 Java student programs---taken from different student populations on the 20-problem institutional corpus. The cross-language C++ cohort ($n=130$) was drawn from an independent semester cohort distinct from the C++ monolingual set ($n=400$); submission ID checks confirmed no overlap between these sets. The pooled-proportion statistic yields $z \approx 1.44$, $p \approx 0.150$ ($\hat{p} = 0.886$, $n_1=n_2=400$). Given the value of p being greater than 0.05, the null hypothesis is accepted: There is no statistical evidence for significant accuracy reduction due to moving from monolingual to cross-language retrieval with the fixed Phase~1 weights.

The optimum for C++ monolingual population differs significantly from the Phase~1 optimum ($w_{BL}=0.00$, $w_{WL}=0.05$, $w_{CES}=0.95$), implying that the Phase~1 weight values represent a stable system behavior point, not a globally optimal configuration per language. Sensitivity analysis reveals that minor fluctuations around the Phase~1 optimum do not impact performance in a significant way (Table~\ref{tab:weight_stability}).

\textbf{Reference pool scale and accuracy context:} The evaluation corpus contains 20 programs per problem, yielding approximately 6--7 cross-language references per query per language. Top-1 retrieval in a candidate pool of this scale is substantially more tractable than any realistically sized deployment scenario. With 6--7 cross-language references per query per language, each retrieval decision is made from a much smaller candidate set than would be encountered in institutional deployments with hundreds of submissions per problem; the reported accuracy figures are operationally valid within this scope and must not be extrapolated to larger reference pools without empirical validation at scale. Characterizing degradation as the reference pool expands to 100+ submissions per problem from multiple institutions remains a priority for future work.

\subsection{Neural Baseline Positioning by Zero-Shot and Fine-Tuned Evaluation}
\label{subsec:neural_comparison}

The system reaches an aggregate cross-language retrieval accuracy of 87.25\%. As neural reference baselines, three models were tested zero-shot with respect to the 400-programs-evaluation set: CodeBERT~\cite{ref_codebert}, which achieved 59.50\% cross-language retrieval; GraphCodeBERT~\cite{ref_graphcodebert}, scoring 58.90\% [95\% confidence interval (CI): 54.01\%, 63.62\%]; and UniXcoder~\cite{ref_unixcoder} scoring 68.92\% [95\% CI: 64.22\%, 73.26\%].

The GraphCodeBERT score (58.90\%) should be considered as a lower bound estimate of its capacity, rather than a representative zero-shot performance estimate: the DFG-conditioned pretraining advantage of GraphCodeBERT demands explicit generation of the DFGs in the preprocessing stage, which was not done here. Tested on plain program token sequences, GraphCodeBERT acts identically to a regular RoBERTa encoder and hence delivers similar results to CodeBERT. This particular experimental setup was chosen because of its transparency as shown in Table~\ref{tab:baseline_comparison}. A proper preprocessed GraphCodeBERT evaluation is left for future work. On the other hand, UniXcoder has been pretrained for cross-lingual code-to-code retrieval without any need for additional graph preprocessing and thus constitutes a more representative neural baseline estimate on this task.

As two fine-tuned neural baselines are used for comparison, CodeBERT fine-tuned (see details in Section~\ref{subsec:codebert_finetuning}) achieves 66.25\% [95\% CI: 55.36\%, 75.65\%] on $n = 80$. This is a +6.75~percentage points improvement over the zero-shot counterpart (59.50\%, $n = 400$). The reason is in the benefit expected from fine-tuning a neural model in the contrastive way on same-strategy cross-language pairs. A limitation of using a three epoch CPU-only training in this case is obvious; a part of the discrepancy with fine-tuned UniXcoder (71.37\%) may be attributed precisely to this factor. Another part is connected with architectural differences between cross-lingual pretraining of different models. Finally, the comparison is possible due to the availability of a fine-tuned UniXcoder instance achieving 71.37\% on the same set (the entire protocol in Section~\ref{subsec:neural_finetuning}). The new proposed framework achieves 88.75\% on the same 80-program held-out test set --- which means 22.50~and 17.38~percentage points higher accuracy than fine-tuned CodeBERT and UniXcoder correspondingly, respectively. An important remark must be made concerning this comparison: both fine-tuned systems have been trained on 320 programs---a significantly smaller amount of the labeled data than is expected for practical application. Certainly with much larger labeled corpora for fine-tuning, the neural methods like UniXcoder will perform better than the suggested approach; the comparison done here assumes the lack of labeled data, i.e., it is focused on the educational deployment scenario. Thus, the deterministic approach that does not require any labeled data provides an obvious deployment advantage under this condition.

\subsection{CES Taxonomy Scalability, Automation, and Extension Complexity}
\label{subsec:ces_stability}

The Context Normalization Protocol uses syntactic expressions as an input, and generates abstract semantic tokens like \texttt{loop\_ANY} or \texttt{rec\_ANY} as its output. In the small world of imperative algorithmic problems that are part of the CS-1 curriculum, the list of canonical computational schemes used to solve such problems – iteration, accumulation, selection, search with early exit, recursion, etc. – is finite and universal, because it is driven by constraints of the particular domain, rather than structural similarities across languages. The CES pattern space is thus strongly aligned in the context of CS-1 programs implemented in C, C++ or Java.

\textbf{Taxonomy scalability.} The 21-pattern taxonomy was developed based on expert examination of the CS-1 corpus, which consists of 20 problems. Every time a new problem space is added to the existing set, (i)~the canonically employed approaches for that particular space need to be found out, (ii)~the CPG rules corresponding to such an approach have to be written, and (iii)~these rules have to be validated on examples which have been left aside during the process. \textcolor{blue}{For a problem in the CS-1 loop-and-recursion space, any extension will be incremental: the P15 failure to recognize the usage of the Java hash set shows how the addition of one pattern can involve a single additional entry in the list of Java API mappings. On the other hand, the number of new patterns required in case of problems using data structures, object-oriented inheritance, or templates becomes super-linear, since in such a situation the vocabulary of context tokens needs to be extended beyond \texttt{loop\_ANY} and \texttt{rec\_ANY} to cover new structural archetypes.} The evaluation corpus (20 problems, 21 patterns) provides strong internal validation but does not characterize this scaling behaviour; multi-institution empirical scaling studies remain future work.

\textbf{Automation possibilities.} The current taxonomy is completely hand-crafted, and its definition requires hand-crafted entries---triples consisting of context, CES name, and predicate. \textcolor{blue}{One obvious way to go towards automation is pattern mining based on CPG-annotated submission corpora---given a labeled set of strategies, frequent subgraph patterns over CPG abstract syntax tree can be discovered as candidate taxonomy entries. This approach will follow the spirit of previous research in vulnerability pattern mining using CPG~\cite{ref_yamaguchi2014}.} Such an effort requires annotated data, a resource which we do not have outside of our single-institution study at the moment. 23-pattern detection rules and STL mapping table with 19 patterns are specific examples when automated pattern frequency discovery would help rule authoring process.

\textbf{Extensions complexity.} The transition from procedural programming in C/C++/Java to dynamically typed languages (Python, JavaScript) adds two more factors of complexity, apart from increasing pattern vocabulary. First, CPG extraction becomes less precise: the Joern's frontend for Python (\texttt{pysrc2cpg}) yields notably less dense sets of nodes compared to \texttt{c2cpg} or \texttt{javasrc2cpg} for the same algorithms; therefore, the chain of normalizations needs to be validated separately per frontend. Second, lack of explicit types in Python/JavaScript complicates CES predicate development, as the \texttt{ACCUMULATIVE::COLLECT} pattern matches Java method \texttt{list.add()} due to the type-resolvability of the API call within the CPG, while its counterpart in Python needs data flow analysis unavailable with the current extraction scripts. This does not mean that extending the tool cannot be done; rather, it implies that every new frontend will require engineering effort proportional to its structural distance from the C-family CPG format.


\subsection{APM Independence and Interpretation}
\label{subsec:apm_discussion}

The APM module and similarity pipeline both operate over the graph representation of CPGs; however, these are architecturally unrelated extraction paths. While the behavioral agreement result generated by the former does not give any ground either for or against the similarity pipeline performance, and vice versa; the two modules are unrelated besides their common graph structure.

The limitation here, however, is one of test suite provenance rather than objective output criteria: the 20 reference source programs were written by the lead author, the test inputs were written by the lead author, and correctness verification was performed using 3-5 structured test cases per each directional pair. The problem here is not the binary correctness criterion --- stdout comparison is objective and deterministic --- but the fact that cases written by the architect of a given system do not necessarily explore the same failure modes as external testers might. In this regard, Problem 8 (GCD) fails when tested on this small test suite precisely because of a condition on the boundary that a third party would be likely to consider. Fuzz testing of the entire APM generative rule base for a larger problem set using externally authored suites is a matter for the future. The takeaway from this result is that this demonstrates the \textit{feasibility}, not necessarily the efficiency bound, of deterministic CPG-based code generation within the restricted CS-1 scope. The archived repository (Joern v1.1.1, JVM 11) reproduces the above-cited 98.3\% figure; the GCD ternary operator limitation is mentioned in the issue tracker for said repository.

\subsection{Threats to Validity}
\label{subsec:threats}

Threats are organized into five thematic groups covering the primary validity dimensions.

\textbf{Threat Group 1 --- Dataset Scope and Annotation Reliability (External and Internal Validity):}
The test set is limited to a subset of 20 problems in CS-1, setting an external validity boundary; further multi-institutional validation is needed to prove external validity. The internal validity relies on the ground truth labels for 400 cross-language submissions being done by the lead author, a possible source of bias, but one that was addressed through the IRR analysis conducted on $n=60$ submissions stratified based on language pairs and difficulty level by an independent domain expert. The interrater agreement achieved is high ($\kappa = 0.86$, $p < 0.001$). An adversarial worst-case scenario, to set a lower bound on our absolute transparency, is calculated when the submissions labeled with uncertainty among raters ($n=93$) were automatically flagged as retrieval failures, which would give our framework a lower limit accuracy of $64.0\%$ ($\frac{349-93}{400}$).

\textbf{Threat Group 2 --- Coverage and Extraction Scope (Construct Validity):}
CES taxonomy includes 21 patterns designed for CS-1 imperative algorithmic problems. Patterns that fall outside this scope---namely, lambda-captured accumulations, iteration via coroutines, C++17/20 range pipeline patterns and Java stream collector patterns---are excluded from the pattern vector. Entry point detection is used for 34.4\% of zero-CES failure in the uncurated ablation dataset, where the program entry point is compared against a curated list of 23 patterns specified in the repository README file. Per problem analysis (Section~\ref{subsec:per_problem}) provides examples of taxonomy gaps: P15 (Java hash-set pattern not included in Java API mapping table) and P14 (variety of reversal-based rotation patterns). If a code submission obtains zero CES scores, 90.3\% accuracy can be obtained with WL fallback out of 93 cases in the ablation set; however, no instruction at the pattern level is provided (interpretability gap). Additionally, the hand-crafted rule-base exhibits inherent brittleness: it is tightly coupled to the specific AST schema of Joern v1.1.1, and minor parser version upgrades or idiosyncratic syntax representations can disrupt pattern matching. Under zero-CES fallback, where submissions employ structural variations not covered by the 21 predicates, the system defaults to WL structural similarity alone, sacrificing semantic interpretability. Scaling this manual rule-authoring process to broader algorithmic domains or object-oriented design patterns is limited by the exponential growth of AST permutation states, contrasting with the automated feature acquisition of deep neural baselines.

\textbf{Threat Group 3 --- Shared Frontend and Directionality of Sparsity (Construct Validity):}
As C and C++ utilize the same frontend (\texttt{c2cpg}) to extract the Joern CPG representation, 246 out of 400 language pairs (61.5\%) are examined with a framework that operates on structurally very similar intermediate representation. The 87.25\% mean includes both the cross-frontend (Java $\rightarrow$ C/C++, $n=140$, 91.43\%) and shared frontend (C $\leftrightarrow$ C++, $n=246$, 86.99\%) results. Directions C $\rightarrow$ Java ($n=11$) and C++ $\rightarrow$ Java ($n=3$) provide statistical data too sparse for inferential validity and have not been included in any aggregate accuracy measurement. Only directional accuracy has been experimentally verified for the Java $\rightarrow$ C and Java $\rightarrow$ C++ cases.

\textbf{Threat Group 4 --- Weight Optimization and Cross-Validation Scope (Internal Validity):}
The Phase~1 weight set was established by grid search on a 104-program C dataset drawn from the institutional corpus. Five-fold cross validation on the C++ monolingual dataset ($91.00\% \pm 3.3\%$) and Java monolingual dataset ($92.00\% \pm 2.92\%$ on the structurally compliant subset) provides clear stability results within languages. No further cross validation has been conducted on the cross-language dataset; however, the two-proportion z-test for weight transfer stability using cross-language testing ($z \approx 1.44$, $p \approx 0.150$, Section~\ref{subsec:weight_stability_discussion}) and the Tversky parameter sensitivity sweep (Table~\ref{tab:tversky_sensitivity}) provide sufficient evidence for cross-language weight transfer stability. Since the Phase~1 framework is currently a manuscript under review, all weight settings ($w_{BL}=0.35$, $w_{WL}=0.40$, $w_{CES}=0.25$) and Phase~1 baseline accuracy figures are reproduced and independently confirmed through sensitivity sweeps in Table~\ref{tab:weight_stability}. Through the sensitivity analysis (Table~\ref{tab:tversky_sensitivity}), it is established that small deviations from the Phase~1 baseline produce no statistical difference in retrieval accuracy, thus disassociating stability of results with adherence to the exact Phase~1 values.

\textbf{Threat Group 5 --- Scope and Replicability of APM Generative Evaluation (Construct Validity and Replicability):}
APM’s scope of generative evaluation is limited to 20 imperative algorithmic problems from CS-1 level. The reference implementations, inputs, and output validations have been done by the first author himself (author conflict of interest declaration). \textcolor{blue}{Consequently, this closed-loop evaluation does not cover non-standard, adversarial, or poorly structured student submissions, and any generalization claims are deferred until third-party fuzz testing with externally authored test suites is performed.} The test cases per algorithm of 3--5 cases do not establish a lower bound for behavioral equivalence; rather, a single recursive edge case was observed for Problem~8, GCD, in this limited set, indicating that further fuzzing might discover more failure cases. \textcolor{blue}{The code property graph (CPG) generation uses Joern v1.1.1 on JVM~11; any changes in the CPG node schemas in newer versions of Joern might change the behavior of the extraction scripts. All scripts, datasets, Joern version, and JVM configuration are publicly archived at \url{https://github.com/Abhinav-C-Das/ckg-multiview-code-similarity} to enable independent reproduction and stress-testing.}
